\documentclass[11pt]{article}
\usepackage[margin=1in]{geometry}
\usepackage[T1]{fontenc}
\usepackage{lmodern,microtype,amsmath,amssymb,amsthm}
\usepackage[hidelinks]{hyperref}
\hypersetup{pdftitle={Report 137 All orders logarithmic asymptotics for A217057},pdfauthor={},pdfsubject={Uniform Banach estimates, regularized coefficients, and inverse thresholds}}
\pdfinfoomitdate=1\pdftrailerid{}\pdfsuppressptexinfo=15
\setlength{\emergencystretch}{2em}
\allowdisplaybreaks[2]
\newtheorem{theorem}{Theorem}[section]
\newtheorem{lemma}[theorem]{Lemma}
\newtheorem{proposition}[theorem]{Proposition}
\newtheorem{corollary}[theorem]{Corollary}
\newcommand{\N}{\mathbb N_0}
\newcommand{\B}{\mathcal B}
\newcommand{\E}{\mathbb E}
\newcommand{\dd}{\,\mathrm d}
\newcommand{\fall}[2]{#1^{\underline{#2}}}
\newcommand{\pa}[2]{\langle #1,#2\rangle_{\mathcal K}}
\title{All orders logarithmic asymptotics for A217057}
\author{Report 137}
\date{October 2, 2026}
\begin{document}
\maketitle
\begin{abstract}
For the number $u_n$ of permutations with exactly one classical occurrence
of $1234$, and its avoidance counterpart $A_n$, we prove a complete
Poincare expansion whose coefficients are affine functions of $\log n$.
The first logarithm occurs at relative order three. Every logarithmic
coefficient is a computable rational multiple of $\sqrt3/\pi$; the next
one after Report 135 is $-5041845\sqrt3/\pi$ at relative order four.
Every nonlogarithmic coefficient has an absolutely convergent regularized
series prescription. In particular, the previously undetermined finite
third coefficient exists. The proof uses uniform all-order character
estimates in a weighted Banach space and a coefficient-level convolution
lemma. No analytic continuation or singularity-analysis hypothesis for
a remainder generating function is assumed. Inverse thresholds have
expansions to every fixed order as well.
\end{abstract}

\section{Exact inputs and the conclusion}
An occurrence of the classical pattern $1234$ is a choice of four
positions with increasing values, without any adjacency requirement.
Let $u_n$ count permutations of $[n]$ with exactly one such choice; set
$u_0=u_1=u_2=u_3=0$. This is OEIS A217057. The corresponding avoidance
count is $A_n$, with $A_0=1$.

Report 134~\cite{R134} established an exact two-half decomposition and a
positive leading ratio. Report 135~\cite{R135} obtained its first two
power corrections and its first logarithmic correction, with an
$O(n^{-3})$ remainder in the ratio. The present report proves the finite
third constant and extends the expansion to every fixed order. All new
uniform estimates and convolution arguments are supplied below. The
unchanged final Report 135 bundle, including its Report 134 companion,
is included with the reproducibility archive for the exact combinatorial
inputs. The advance claimed here is relative to those reports; an
exhaustive literature-priority claim is not made.

Let $F_m(p,q)$ count $1234$-avoiders with the final $p$ entries decreasing
and with the $q$ smallest values occurring in decreasing order. Write
$A_m=F_m(0,0)$, and put
\[
H_s(i,j)=F_{s+2}(i+1,j+1)-F_{s+1}(i,j)\quad(i+j\le s),
\qquad H_s(i,j)=0\quad(i+j>s).
\]
For $0\le p,q\le m$, we use the following proved inputs from
Report 135, whose combinatorial foundations are supplied in its
Report 134 companion:
\begin{align}
u_n&=\sum_{s+t=n-4}\pa{H_s}{H_t},\label{eq:glue}\\
\pa{X}{Y}&=\sum_{i,j,k,l\ge0}X(i,j)Y(k,l)
 \binom{i+k}{i}\binom{j+l}{j},\label{eq:pair}\\
0\le H_s(i,j)&\le K A_s(i+1)^4(j+1)^4 3^{-i-j},\label{eq:bound}\\
c\frac{9^m}{(m+1)^4}&\le A_m\le C\frac{9^m}{(m+1)^4},\label{eq:Abound}\\
F_m(p,q)&=\int_{U(3)}P^{m-p}h_p\overline{P^{m-q}h_q}\dd U.
\label{eq:unitary}
\end{align}
Here $P=\operatorname{tr}U$, $h_p$ is the complete homogeneous character,
and Haar measure has mass one. Set
\[
C_0=\frac{81\sqrt3}{16\pi},\quad d_p=\binom{p+2}{2}3^{-p},
\quad a_p=\frac{p(p-1)}2,
\]
\begin{align}
C(i,j)&=81d_{i+1}d_{j+1}-9d_id_j,\label{eq:C}\\
B(i,j)&=-81d_{i+1}d_{j+1}(8+a_{i+1}+a_{j+1})
 +9d_id_j(4+a_i+a_j).\label{eq:B}
\end{align}
The avoidance and half expansions of Report 135 start as
\[
A_m=C_0 9^m m^{-4}\left(1-\frac{11}{2m}+\frac{20}{m^2}
 +O(m^{-3})\right),\qquad
\frac{H_m(i,j)}{A_m}=C(i,j)+\frac{B(i,j)}m+\cdots.
\]
The positivity $R>0$ of its ratio limit is another inherited input.

\begin{theorem}\label{thm:main}
There are uniquely determined real numbers $r_j,\kappa_j$ such that,
for every fixed integer $J\ge0$,
\begin{equation}\label{eq:full}
\frac{u_n}{A_n}=\sum_{j=0}^J\frac{r_j+\kappa_j\log n}{n^j}
 +O_J\left(\frac{1+\log n}{n^{J+1}}\right).
\end{equation}
They satisfy
\[
r_0=R,\quad r_1=S,\quad r_2=S_2,\qquad
\kappa_0=\kappa_1=\kappa_2=0,
\]
\begin{equation}\label{eq:kappas}
\kappa_3=\frac{43020\sqrt3}{\pi},\qquad
\kappa_4=-\frac{5041845\sqrt3}{\pi}.
\end{equation}
Every $\kappa_j$ belongs to $(\sqrt3/\pi)\mathbb Q$ and is computable
by finite rational operations. Each $r_j$ is specified by finitely many
absolutely convergent Banach-valued regularized series followed by
continuous pairings and finite algebra with universal harmonic constants.
In particular,
\[
\lim_{n\to\infty}\left[n^3\left(\frac{u_n}{A_n}
 -R-\frac Sn-\frac{S_2}{n^2}\right)-\kappa_3\log n\right]=r_3
\]
exists; an explicit convergent-series recipe is given in
Section~\ref{sec:third}.
\end{theorem}
This is an all-orders asymptotic statement at each fixed order. It makes
no assertion of convergence or summability of the infinite asymptotic
series, or of an expansion into exponentially small sectors. It supplies
neither certified decimal values for the regularized constants nor an
effective numerical onset.

\section{All orders estimates in a weighted Banach space}
Fix $1/3<r<1/2$, and let $E_r$ be the space of real arrays with norm
\[
\|X\|_r=\sup_{i,j\ge0}|X(i,j)|r^{-i-j}.
\]
It is a Banach space, by its isometric identification with bounded real
arrays after multiplying coordinate $(i,j)$ by $r^{-i-j}$. The pairing in~\eqref{eq:pair} is symmetric and
continuous, since
\begin{equation}\label{eq:pairnorm}
|\pa XY|\le(1-2r)^{-2}\|X\|_r\|Y\|_r.
\end{equation}
Indeed $\sum_{i,k\ge0}r^{i+k}\binom{i+k}{i}=(1-2r)^{-1}$.
Put $a_n=9^{-n}H_n\in E_r$. The global estimates imply
$\|a_n\|_r=O((n+1)^{-4})$, since a fixed polynomial times
$(3r)^{-i-j}$ is bounded. We reserve $A_m$ for avoidance counts and use $U_j$ for the Banach
coefficient arrays below.

\begin{proposition}\label{prop:Banach}
There are arrays $U_j\in E_r$, independent of $r$, such that for every
fixed $M\ge0$,
\begin{equation}\label{eq:Banach}
\left\|a_n-n^{-4}\sum_{j=0}^M U_j n^{-j}\right\|_r
 =O_{M,r}(n^{-M-5}).
\end{equation}
They have the form $U_j(i,k)=C_0 3^{-i-k}P_j(i,k)$ for a rational
polynomial $P_j$. The first two are
\begin{equation}\label{eq:A01}
U_0=C_0C,\qquad U_1=C_0\left(-\frac{11}{2}C+B\right).
\end{equation}
The avoidance sequence has a complete pure-power expansion
\begin{equation}\label{eq:avoidall}
A_m=C_0 9^m m^{-4}\left(\sum_{j=0}^M\alpha_jm^{-j}
 +O_M(m^{-M-1})\right),
\quad\alpha_j\in\mathbb Q,
\end{equation}
where $\alpha_0=1$, $\alpha_1=-11/2$, $\alpha_2=20$, and
$\alpha_3=-965/16$.
\end{proposition}

\subsection{The fixed order Laplace estimates}
Common central phase cancels in~\eqref{eq:unitary}. On the quotient
of the eigenangle torus by common phase, $|P|=3$ has a single maximum.
Use a small centered chart $\theta_1+\theta_2+\theta_3=0$, and put
$Q=\sum\theta_i^2$ and $\Delta=\prod_{i<j}(\theta_i-\theta_j)$.
The phase and factored Weyl density are real analytic there:
\[
\phi(\theta)=\log(|P|^2/9)=-Q/3+O(Q^2),\qquad
\prod_{i<j}|e^{i\theta_i}-e^{i\theta_j}|^2
 =\Delta(\theta)^2 w(\theta),\quad w(0)=1.
\]
Both $\phi$ and $w$ are even under $\theta\mapsto-\theta$.
In a sufficiently small chart $\phi\le-cQ$; outside it the trace ratio
is at most some $\rho<1$. The quotient has a constant local Jacobian.
Rescaling $\theta=y/\sqrt m$ gives leading density
$e^{-Q_y/3}\Delta(y)^2\dd y$ and a leading factor $m^{-4}$.

Here is a fixed-order remainder argument. Write
$\phi=-Q/3+v$, where $v=O(Q^2)$, and fix an integer $M\ge0$.
Choose the chart small enough that $|v|\le Q/12$. Truncate
$v$ through degree $2M+2$, leaving a remainder $O(Q^{M+2})$,
and truncate $w$ through degree $2M$, leaving $O(Q^{M+1})$.
For $M=0$ the first truncation is the empty polynomial.
Changing the phase to that polynomial creates an error bounded by
$C mQ^{M+2}e^{-c'mQ}$, by the mean value formula for the exponential;
its integral is of relative order $m^{-M-1}$.

For the polynomial phase perturbation $v_{\le2M+2}$, Taylor's integral
formula bounds the remainder after its exponential is expanded through
power $M$ by
\[
C_M m^{M+1}Q^{2M+2}e^{m|v_{\le2M+2}|}.
\]
A smaller chart absorbs this exponential into $e^{-mQ/3}$, and its
integral has relative order $m^{-M-1}$. Expand the retained finite
polynomial products and retain exactly the monomials of relative
weight at most $M$, where $m^a$ times a homogeneous polynomial of
degree $2b$ has weight $b-a$. Each discarded term has weight at least
$M+1$. In fact, direct radial scaling on the two-dimensional plane gives
\[
\int Q^b\Delta^2 e^{-c'mQ}\dd\theta=O(m^{-4-b});
\]
thus $m^aQ^b$ has the claimed relative order. The amplitude remainder is
handled by the same bound. Extending the resulting Gaussian polynomial
integrals to the whole plane changes them by an exponentially small
quantity. Odd terms integrate to zero on an inversion-symmetric chart.
This proves the fixed-order expansion with an actual remainder, rather
than only a formal series.

The argument also applies with an arbitrary fixed polynomial insertion.
For a homogeneous insertion of degree $2h$, scaling first extracts
$m^{-h}$, and the proof gives any requested number of further integer
powers. Linearity and finite-dimensional norm equivalence make the
remainder constant uniform on a coefficient-norm unit ball of any
fixed polynomial space. The complementary region is exponentially
small, using the original integrand there.
The same argument gives, for fixed $c_1,h$ and sufficiently small
$\varepsilon>0$,
\begin{equation}\label{eq:moment}
\E_m[Q^h e^{c_1LQ};\mathrm{local}]\le C_hm^{-h}
\quad(0\le L\le\varepsilon m),
\end{equation}
where $\E_m$ is normalized trace-moment expectation.
These arguments prove~\eqref{eq:avoidall} and all polynomial-insertion
moment expansions, uniformly on the unit ball of each fixed
finite-dimensional polynomial space.

The normalized Gaussian moments are rational. For example, write
$\theta=(x,y,-x-y)$: the Gaussian $e^{-Q/3}$ has covariance
$\left(\begin{smallmatrix}1&-1/2\\-1/2&1\end{smallmatrix}\right)$,
and the additional weight is the rational polynomial $\Delta^2$.
The Gaussian moment recurrence of Isserlis~\cite{Isserlis}, often
called Wick's rule, then gives rational normalized moments. All Taylor
coefficients above are rational, proving $\alpha_j\in\mathbb Q$.
The constant $C_0$ is the already established leading constant.

\subsection{Uniformity in both boundary parameters}
Put $c_p=\binom{p+2}{2}$ and
\[
r_p(\theta)=\frac{h_p(e^{i\theta})}{c_p(P(\theta)/3)^p},
\qquad h_p/P^p=d_pr_p.
\]
For a centered unit vector $v$, set $g_p(t)=h_p(e^{itv})/c_p$ and
$z(t)=P(tv)/3$. For every fixed derivative order $D$,
\[
|g_p^{(h)}(t)|\le(p+1)^h\quad(0\le h\le D),\qquad
\left|\frac{\mathrm d^h}{\mathrm dt^h}z(t)^{-p}\right|
 \le C_D(p+1)^h e^{C_Dpt^2}
\]
on a fixed sufficiently small interval, uniformly in $v$. The first
bound follows by averaging the $c_p$ weak-composition terms; the second
is the finite chain rule, using $|z(t)|\ge e^{-Ct^2}$.
Consequently, for $L=p+q$ and any fixed $M$, the real even function
$\Re(r_p\overline{r_q})$ has a homogeneous Taylor expansion
\begin{equation}\label{eq:boundaryTaylor}
\Re(r_p\overline{r_q})=\sum_{h=0}^M D_{2h}^{p,q}(\theta)
 +\mathcal R_{M,p,q}(\theta),
\end{equation}
with
\[
|D_{2h}^{p,q}(\theta)|\le C_M(1+L)^{2h}Q^h,
\quad
|\mathcal R_{M,p,q}(\theta)|
 \le C_M(1+L)^{2M+2}Q^{M+1}e^{C_MLQ}.
\]
The real product is even because the characters conjugate when the
angles are negated. Taylor's theorem through degree $2M+1$ proves the
remainder bound; its last odd homogeneous coefficient vanishes.

Integrate the remainder using~\eqref{eq:moment}. Expand the expectation
of $D_{2h}^{p,q}$ through relative order $M-h$ by the polynomial-insertion
argument above. Uniformity on that finite-dimensional space, and the
displayed polynomial bound for its coefficients, show that for some
$\varepsilon_M>0$ and uniformly when $p+q\le\varepsilon_M m$,
\begin{equation}\label{eq:Fhigher}
\frac{F_m(p,q)}{A_m}=d_pd_q\left[
 \sum_{j=0}^M f_j(p,q)m^{-j}
 +O_M\left((1+p+q)^{2M+2}m^{-M-1}\right)\right].
\end{equation}
The $f_j$ are rational polynomials of degree at most $2j$, with $f_0=1$.
To see polynomiality directly, weak compositions satisfy
\[
\frac1{c_p}\sum_{\delta_1+\delta_2+\delta_3=p}
 (\delta_1)_{\underline a}(\delta_2)_{\underline b}
 (\delta_3)_{\underline c}
 =\frac{2a!b!c!}{(a+b+c+2)!}\fall p{a+b+c}.
\]
All Taylor coefficients of the averaged character, and of $z^{-p}$,
are therefore rational polynomials in the parameters.

There is no division by $P$ outside the chart. The original polynomial
integrand in~\eqref{eq:unitary} has modulus bounded by
\[
9^m d_pd_q\left(\frac{|P|}{3}\right)^{2m-p-q}.
\]
After division by $A_m$ its off-chart part is
$O(d_pd_qm^4\rho^{(2-\varepsilon_M)m})$, absorbed in the remainder.
Thus~\eqref{eq:Fhigher} is uniform on its entire stated range.

Combining~\eqref{eq:Fhigher},~\eqref{eq:avoidall} and the fixed shifts
$m=n+1,n+2$ yields the claimed half coefficients and, whenever
$i+j+2\le\varepsilon'_M n$, a bound
\[
\left|a_n(i,j)-n^{-4}\sum_{h=0}^M U_h(i,j)n^{-h}\right|
 \le C_Mn^{-M-5}(1+i+j)^{D_M}3^{-i-j}
\]
for a fixed finite $D_M$. Dividing by $r^{i+j}$ makes the boundary
polynomial uniformly bounded. In the complementary range,
\eqref{eq:bound} and the polynomial coefficient bounds give an
exponentially small norm tail, since $(3r)^{-i-j}$ decays geometrically
and $i+j\ge\varepsilon'_Mn-2$. This includes the triangular zero region
$i+j>n$. It proves~\eqref{eq:Banach} and completes the uniform estimates.


\subsection{A finite algorithm for every local coefficient}\label{sec:algorithm}
The construction above can be carried out with rational polynomials
at each fixed order. Let $\phi_{2h}$ and $w_{2h}$ be the homogeneous
parts of the local phase and amplitude, where $\phi_2=-Q/3$ and
$w_0=1$. With $t$ a formal variable, define
\begin{equation}\label{eq:Kalgorithm}
K_j(y)=[t^j]\left\{
 \left(\sum_{h=0}^j t^h w_{2h}(y)\right)
 \exp\left(\sum_{h=2}^{j+1}t^{h-1}\phi_{2h}(y)\right)
\right\}.
\end{equation}
Only finitely many terms of the exponential can contribute. If
$\E_0$ is expectation for the normalized density
$e^{-Q/3}\Delta^2$, then $\alpha_j=\E_0K_j$. The normalization agrees
with $\alpha_0=1$; the leading factor $C_0$ is the inherited one.

For the boundary insertion in~\eqref{eq:boundaryTaylor}, put
\[
N_j(p,q)=\sum_{h=0}^j\E_0\bigl[D_{2h}^{p,q}K_{j-h}\bigr].
\]
The finite division rule
\begin{equation}\label{eq:falgorithm}
f_0=1,\qquad
f_j=N_j-\sum_{h=1}^j\alpha_h f_{j-h}
\end{equation}
gives the coefficients in~\eqref{eq:Fhigher}. The averaged character
Taylor coefficients are obtained from the displayed weak-composition
factorial moments, converting ordinary powers to falling powers if
necessary. The reciprocal trace is expanded by its finite binomial
series. Thus this is an explicit rational algorithm for the polynomials
$f_j$, including $f_2$ and $f_3$, without needing a separate closed-form
third-boundary formula.

For clarity, it also gives the half arrays directly. Write
$g_j(p,q)=\sum_{h=0}^j\alpha_h f_{j-h}(p,q)$, and define, for a fixed
integer shift $s$,
\[
E_j(s;p,q)=\sum_{h=0}^j g_h(p,q)
 \binom{-4-h}{j-h}s^{j-h}.
\]
Expanding $(n+s)^{-4-h}$ proves
\begin{equation}\label{eq:Aalgorithm}
\begin{split}
U_j(i,k)=C_0\bigl[&81d_{i+1}d_{k+1}E_j(2;i+1,k+1)\\
 &-9d_id_kE_j(1;i,k)\bigr].
\end{split}
\end{equation}
These formulas use only finitely many exact Gaussian moments at any
fixed $j$. One elementary implementation, in coordinates $(x,y,-x-y)$,
sets $G(a,b)=\E_{\rm G}(x^ay^b)$ and uses
\[
G(a,b)=(a-1)G(a-2,b)-\frac b2G(a-1,b-1)\quad(a\ge1),
\]
with $G(0,2b)=(2b-1)!!$, $G(0,2b+1)=0$, $G(0,0)=1$, and negative
indices interpreted as zero. This recurrence follows directly by
Gaussian integration by parts. Multiplying by $\Delta^2$ and dividing
by $\E_{\rm G}\Delta^2$ then evaluates $\E_0$.
\subsection{A reproducible check of the third avoidance coefficient}
Let $D_1=x-y,D_2=x+2y,D_3=2x+y$, and let $L_{2h}$ be the coefficient
of $t^h$ in
\[
\log\left(1+\frac29\sum_{h\ge1}
 \frac{(-1)^ht^h}{(2h)!}\sum_{i=1}^3D_i^{2h}\right).
\]
Let $W_{2h}$ be the coefficient of $t^h$ in
\[
\prod_{i=1}^3\left(\frac{\sin(\sqrt tD_i/2)}{\sqrt tD_i/2}\right)^2.
\]
Expectation below is with Gaussian covariance just specified and extra
weight $\Delta^2=D_1^2D_2^2D_3^2$. Its normalizing Gaussian moment is
$\E_{\rm G}\Delta^2=81/2$. Then
\begin{align*}
\alpha_3=\E_0\bigl[W_6+W_4L_4+W_2(L_6+L_4^2/2)
 +L_8+L_4L_6+L_4^3/6\bigr]=-965/16.
\end{align*}
These are finite polynomial calculations. The companion replay evaluates the Gaussian moments exactly over
$\mathbb Q$, reproducing $1,-11/2,20,-965/16$.

\section{A coefficient level convolution theorem}
The following elementary argument works for any real Banach space $E$
with a continuous symmetric bilinear form $\B:E\times E\to\mathbb R$.
Assume a sequence $a_n\in E$ has~\eqref{eq:Banach}, with the same leading
power $n^{-4}$. Let
\[
v_n=\sum_{t=0}^n\B(a_t,a_{n-t}).
\]
For $k\ge3$ define the scalar model sequence
\begin{equation}\label{eq:bk}
b_k(n)=[z^n](1-z)^k\log(1-z).
\end{equation}
For $n>k$ the coefficient is exactly
\begin{equation}\label{eq:bexact}
b_k(n)=\frac{(-1)^{k+1}k!}{\fall n{k+1}}.
\end{equation}
For all $n\ge0$ the coefficient is also given by the finite sum
\[
b_k(n)=-\sum_{h=0}^{\min(k,n-1)}\frac{(-1)^h\binom kh}{n-h},
\qquad b_k(0)=0,
\]
with the sum used when $n\ge1$. For example,
$(b_3(0),b_3(1),b_3(2),b_3(3))=(0,-1,5/2,-11/6)$.
The finitely many coefficients at $n\le k$ retain their definition
in~\eqref{eq:bk}; replacing them by another extension would change the
regularized constants and is not permitted.

\begin{lemma}\label{lem:models}
There are unique vectors $V_k\in E$, $k\ge3$, such that for every
$K\ge3$,
\begin{equation}\label{eq:rdecay}
r^{(K)}_n=a_n-\sum_{k=3}^KV_kb_k(n)=O_E((n+1)^{-K-2}).
\end{equation}
If $h_d(1,\ldots,k)$ denotes the complete homogeneous polynomial of
degree $d$ in $1,\ldots,k$, they are computed by
\begin{equation}\label{eq:Vrec}
V_k=\frac{U_{k-3}-\sum_{\ell=3}^{k-1}
 (-1)^{\ell+1}\ell!h_{k-\ell}(1,\ldots,\ell)V_\ell}
 {(-1)^{k+1}k!}.
\end{equation}
\end{lemma}
\begin{proof}
Expand the rational function~\eqref{eq:bexact} as
$(-1)^{k+1}k!n^{-k-1}\prod_{j=1}^k(1-j/n)^{-1}$ to any fixed order.
Its leading coefficient is nonzero. Triangular matching against the
assumed Banach expansion gives~\eqref{eq:Vrec} and~\eqref{eq:rdecay}.
\end{proof}

For $j\ge0$ define a regularized Taylor coefficient by
\begin{equation}\label{eq:Tj}
T_j=\frac{(-1)^j}{j!}\sum_{n\ge0}\fall nj
 \left(a_n-\sum_{k=3}^{j}V_kb_k(n)\right).
\end{equation}
An empty inner sum means zero. These series converge absolutely in $E$:
for $j\le2$, use $a_n=O_E(n^{-4})$; for $j\ge3$, apply
\eqref{eq:rdecay} with $K=j$. In the latter case the weighted summand is
$O_E(n^{-2})$.

\begin{lemma}\label{lem:remainder}
For a fixed $K\ge3$ set
\[
S_K(z)=\sum_{k=3}^KV_k(1-z)^k\log(1-z),\qquad
P_K(z)=\sum_{j=0}^{K-1}T_j(1-z)^j.
\]
Then, as $n\to\infty$,
\begin{equation}\label{eq:finitecalc}
v_n=[z^n]\left(2\B(P_K,S_K)+\B(S_K,S_K)\right)
 +O(n^{-K-2}).
\end{equation}
\end{lemma}
\begin{proof}
Let $R_K(z)=\sum r_n^{(K)}z^n$. Its derivatives through order $K$ at
$z=1$ are absolutely convergent coefficient sums. For $k>j$, the
$j$th derivative of $(1-z)^k\log(1-z)$ tends to zero at $z=1$ and its
coefficient derivative series is absolutely convergent. Therefore the
Taylor polynomial of $R_K$ through degree $K-1$, expressed in powers
of $1-z$, is exactly $P_K$ from~\eqref{eq:Tj}.
Let $q_n=[z^n](R_K-P_K)$. Then
\begin{equation}\label{eq:q}
\|q_n\|=O((n+1)^{-K-2}),\quad
\sum_{n\ge0}\fall njq_n=0\ (0\le j\le K-1),\quad
\sum_{n\ge0}n^K\|q_n\|<\infty.
\end{equation}
The zero falling-factorial moments also imply zero ordinary moments
through $K-1$.

All generating functions converge absolutely for $|z|<1$, so their
products equal coefficient convolutions. Terms involving $P_K$ twice
are polynomial and vanish for sufficiently large $n$. Convolution of
$q$ with $P_K$ is $O(n^{-K-2})$. Convolution of $q$ with itself has the
same bound: split at $n/2$, bound the larger index by $O(n^{-K-2})$,
and sum the other sequence absolutely.

It remains to bound $q*b_k$, for fixed $3\le k\le K$. In the part
$t>n/2$, $\|q_t\|=O(n^{-K-2})$ and $\sum_s|b_k(s)|<\infty$.
In the part $t\le n/2$, use the rational extension in~\eqref{eq:bexact}:
its $h$th derivative on $[n/2,n]$ is $O(n^{-k-1-h})$ for all large $n$.
Taylor expansion of $b_k(n-t)$ through degree $K-1$ has remainder
$O(t^Kn^{-k-1-K})$. Its convolution with $q$ is bounded using the last
condition in~\eqref{eq:q}. Each retained moment, after extending the
sum to all $t$, vanishes. The truncation tail for its degree $h$ is
\[
O(n^{-k-1-h})\sum_{t>n/2}t^h\|q_t\|
 =O(n^{-k-K-2}).
\]
Thus the entire convolution is $O(n^{-K-2})$. Continuity of $\B$ handles
its vector coefficients. Expansion of
$\B(P_K+S_K+(R_K-P_K),P_K+S_K+(R_K-P_K))$ proves the lemma.
No continuation of $R_K$ across $|z|=1$ has been used.
\end{proof}

\begin{theorem}\label{thm:conv}
The convolution has a complete expansion
\begin{equation}\label{eq:convexp}
v_n=\sum_{j=0}^J(c_j+\ell_j\log n)n^{-4-j}
 +O_J(n^{-5-J}(1+\log n))
\end{equation}
for every fixed $J$, with $\ell_0=\ell_1=\ell_2=0$.
Every coefficient is obtained by expanding the finite expression
in~\eqref{eq:finitecalc}, with any $K\ge J+4$.
In particular, each $\ell_j$ depends only on the local tail vectors
$U_h$, never on the regularized moments $T_h$.
\end{theorem}
\begin{proof}
Products $P_KS_K$ are finite sums of the single-logarithm models
in~\eqref{eq:bk}. Products $S_K^2$ are finite sums of
$(1-z)^m\log^2(1-z)$ with $m\ge6$. For $n>m$, differentiation twice
in $\alpha$ of the finite polynomial
$[z^n](1-z)^\alpha=(-1)^n\fall\alpha n/n!$ gives the exact formula
\begin{equation}\label{eq:log2}
[z^n](1-z)^m\log^2(1-z)
 =2b_m(n)(H_m-H_{n-m-1}).
\end{equation}
Here $H_m=\sum_{h=1}^m1/h$ and $H_0=0$; these scalar harmonic
numbers are distinct from the earlier half arrays $H_s(i,j)$.
For a fixed integer $p\ge0$, Euler--Maclaurin gives
\[
H_N=\log N+\gamma+\frac1{2N}
 -\sum_{h=1}^p\frac{B_{2h}}{2hN^{2h}}+O_p(N^{-2p-2}),
\]
where $B_{2h}$ are Bernoulli numbers. The finite remainder follows from
the integral remainder in Euler--Maclaurin applied to $1/x$; the
periodic Bernoulli function is bounded and the appropriate derivatives
are integrable on $[N,\infty)$. Equivalently, see the exact
harmonic--digamma identity and its fixed-order expansion in
DLMF~\cite{DLMFharm,DLMFEM}. Expanding $N=n-m-1$ with fixed $m$ and
using~\eqref{eq:bexact} gives full expansions involving only a constant
and a single $\log n$. The earliest log-squared model has $m=6$, so its
coefficient logarithm first occurs at $n^{-7}$. No higher logarithmic
powers occur. Taking $K\ge J+4$ makes the error of
\eqref{eq:finitecalc} smaller than the error in~\eqref{eq:convexp}.
\end{proof}

\section{The constants and the first two logarithms}\label{sec:third}
Apply Theorem~\ref{thm:conv} to $E_r$, $\B=\langle\cdot,\cdot\rangle_{
\mathcal K}$ and $a_n=9^{-n}H_n$. By~\eqref{eq:glue},
\begin{equation}\label{eq:unshift}
u_n=9^{n-4}v_{n-4}.
\end{equation}
The fixed shift, the pure-power avoidance expansion, and ordinary
formal division yield~\eqref{eq:full}, with its claimed error.
The coefficients are unique: at the first differing order, multiplying
the difference by that power of $n$ would give a nonzero affine
function of $\log n$ tending to zero. Thus the first three constants
agree with $R,S,S_2$ from Report 135.

Here is an explicit recipe for the finite third coefficient. The first
four singular vectors are
\begin{align}
V_3&=U_0/6,&V_4&=(6U_0-U_1)/24,\notag\\
V_5&=(35U_0-10U_1+U_2)/120,&
V_6&=(225U_0-85U_1+15U_2-U_3)/720.\label{eq:Vfirst}
\end{align}
Compute $U_2,U_3$ by the finite character/Gaussian algorithm in
Section~\ref{sec:algorithm}. Form the absolutely convergent series
\begin{align}
T_0&=\sum_{n\ge0}a_n,&T_1&=-\sum_{n\ge0}na_n,\notag\\
T_2&=\frac12\sum_{n\ge0}n(n-1)a_n,&
T_3&=-\frac16\sum_{n\ge0}n(n-1)(n-2)(a_n-V_3b_3(n)).
\label{eq:Tfirst}
\end{align}
For $m=3,4,5,6$, put
\[
\Phi_m=2\sum_{k=3}^{m}\pa{T_{m-k}}{V_k}.
\]
The finite coefficients of $v_n$ through the first log are
\begin{align}
c_0&=6\Phi_3,\notag\\
c_1&=36\Phi_3-24\Phi_4,\notag\\
c_2&=150\Phi_3-240\Phi_4+120\Phi_5,\notag\\
\ell_3&=1440\pa{V_3}{V_3}=40\pa{U_0}{U_0},\notag\\
c_3&=540\Phi_3-1560\Phi_4+1800\Phi_5-720\Phi_6
 +\ell_3(\gamma-H_6),\label{eq:cfirst}
\end{align}
where $\gamma$ is Euler's constant. These formulas use the finite
coefficients of $b_k$ at small $n$ exactly as in~\eqref{eq:bk}.
The inverse avoidance series starts
\[
(1-11/(2n)+20/n^2-965/(16n^3)+\cdots)^{-1}
 =1+\frac{11}{2n}+\frac{41}{4n^2}+\frac{107}{16n^3}+\cdots.
\]
Expanding the shift $n-4$ in~\eqref{eq:unshift} now gives
\begin{equation}\label{eq:r3}
r_3=\frac{9^{-4}}{C_0}
 \left(c_3+\frac{59}{2}c_2+\frac{1441}{4}c_1
 +\frac{37291}{16}c_0\right).
\end{equation}
Equations~\eqref{eq:Tfirst}--\eqref{eq:r3}, together with finite local
polynomial calculations, are the promised absolutely convergent
regularized-series prescription. Existence and convergence, not a
numerically certified evaluation, are claimed here.

The next logarithmic coefficient of the unshifted convolution is
\begin{align*}
\ell_4&=30240\pa{V_3}{V_3}-20160\pa{V_3}{V_4}
       =140\pa{U_0}{U_1}.
\end{align*}
The shift and division therefore give
\begin{align}
\kappa_3&=\frac{9^{-4}}{C_0}\ell_3,
\notag\\
\kappa_4&=\frac{9^{-4}}{C_0}\left(\ell_4+\frac{67}{2}\ell_3\right)
 =C_0 9^{-4}\left(570\pa CC+140\pa CB\right).\label{eq:kappa4}
\end{align}
For transparency, both needed pairings are exact rational calculations:
\begin{equation}\label{eq:pairvals}
\pa CC=1393848,\qquad \pa CB=-52348032.
\end{equation}
A convenient general evaluation rule is this: if
$f_i=3^{-i}f(i)$ and $g_i=3^{-i}g(i)$ for rational polynomials, write
$f(i)=\sum_p f_p\fall ip$ and $g(i)=\sum_q g_q\fall iq$. Then
\begin{equation}\label{eq:fallpair}
\sum_{i,k\ge0} f_i g_k\binom{i+k}{i}
 =3\sum_{p,q}f_pg_q(p+q)!.
\end{equation}
It follows by differentiating $(1-x-y)^{-1}$ and putting $x=y=1/3$.
Absolute convergence licenses every differentiation. Rank-one
factorizations of $C,B$ reduce~\eqref{eq:pairvals} to this rule; the companion replay carries out the rational arithmetic. Substitution proves~\eqref{eq:kappas}.

More generally, every $V_k$ is $C_0$ times an array of the form
$3^{-i-j}$ times a rational polynomial. Thus all pairings of singular
vectors are $C_0^2$ times rationals by~\eqref{eq:fallpair}.
The log coefficient of~\eqref{eq:log2} is the rational function
$-2b_m(n)$. Shift and division use rational coefficients. Consequently
every $\kappa_j$ is a rational multiple of $C_0$, proving the final
algebraic assertion of Theorem~\ref{thm:main}.

\section{All orders inverse thresholds}
Let $M_0=C_0R>0$, $\lambda=\log9$, and
\[
N(T)=\min\{n\ge0:u_n\ge T\}.
\]
The ratio $u_{n+1}/u_n\to9$ makes the sequence eventually strictly
increasing. For sufficiently large thresholds, this controls the
inverse irrespective of finitely many early values.

The direct expansion implies, to every fixed order,
\begin{equation}\label{eq:logu}
\log u_n=\lambda n-4\log n+\log M_0
 +\sum_{j=1}^{J}\frac{L_j(\log n)}{n^j}
 +O_J\left(\frac{(1+\log n)^{J+1}}{n^{J+1}}\right),
\end{equation}
where $L_j$ is a polynomial with $\deg L_j\le\lfloor j/3\rfloor$.
Indeed every logarithm in the normalized direct expansion costs at
least three powers of $1/n$; expand the ordinary scalar logarithm by
its finite Taylor formula. This degree bound concerns $\log u_n$,
where higher powers of logarithms may occur; it does not contradict the
affine-log structure of $u_n/A_n$.

\begin{corollary}\label{cor:inverse}
For every fixed $J\ge0$ there are explicitly recursive polynomials
$Q_j$, $1\le j\le J$, of degree at most $j$, such that, with
$y=\log T$ and $d=\log M_0+4\log\lambda$,
\[
x_J(T)=\frac1\lambda\left[y+4\log y-d
 +\sum_{j=1}^J\frac{Q_j(\log y)}{y^j}\right]
\]
satisfies the eventual sandwich
\begin{equation}\label{eq:inverse}
\left\lceil x_J(T)-C_J\frac{(1+\log y)^{J+1}}{y^{J+1}}\right\rceil
 \le N(T)\le
\left\lceil x_J(T)+C_J\frac{(1+\log y)^{J+1}}{y^{J+1}}\right\rceil.
\end{equation}
The constants and onset are not asserted effective. No convention of
rounding a truncated asymptotic expansion to a single exact integer is
implied when its error interval crosses an integer.
\end{corollary}
\begin{proof}
Set $w=\lambda x$ and formally substitute
$w=y+4\log y-d+\sum_{j=1}^JQ_j(\log y)y^{-j}$ in
\eqref{eq:logu}. At each new order $y^{-j}$ the new polynomial $Q_j$
appears with coefficient one; all other terms have already been
computed. This determines $Q_j$ successively. Expansion of
$\log(1+(w-y)/y)$ and $(1+(w-y)/y)^{-h}$ shows inductively that its
degree is at most $j$. Finite Taylor remainders on $(w-y)/y=o(1)$ give
residual $O_J((1+\log y)^{J+1}y^{-J-1})$.

For a rigorous threshold conclusion, regard the right side of the
finite expansion~\eqref{eq:logu} as a smooth function of a positive
real variable $x$. Its derivative is $\lambda+o(1)$, bounded below by
$\lambda/2$ for all large $x$. The discrete remainder and the
substitution residual have the stated bound uniformly for integers
within a bounded distance of $x_J$. Increasing the constant $C_J$ makes
the lower and upper shifts dominate both errors. For full precision about rounding, put
$\epsilon=C_J(1+\log y)^{J+1}y^{-J-1}$. The lower comparison is at
$k_-=\lceil x_J-\epsilon\rceil-1<x_J-\epsilon$, and the upper one
is at $k_+=\lceil x_J+\epsilon\rceil\ge x_J+\epsilon$.
The derivative bound and the chosen error margin give
$\log u_{k_-}<\log T$ and $\log u_{k_+}\ge\log T$.
All these integers lie within a bounded distance of $x_J$.
Eventual monotonicity, with the finite early maximum below $T$, gives
$k_-<N(T)\le k_+$, exactly~\eqref{eq:inverse}.
\end{proof}
For example, write
$u_n=M_0 9^n n^{-4}(1+\beta/n+\delta/n^2+\cdots)$, where
\[
\beta=-\frac{11}{2}+\frac SR,
\qquad \delta=20-\frac{11S}{2R}+\frac{S_2}{R}.
\]
If
$a=4z-d$, the first two inverse polynomials are
\[
Q_1(z)=4a-\lambda\beta,
\qquad Q_2(z)=4Q_1(z)-2a^2+\lambda\beta a
 -\lambda^2(\delta-\beta^2/2).
\]
At higher orders the same recursion includes the logarithmic
coefficients and the regularized constants from the direct expansion.

\section{Reproducibility and the scope of the computations}
The companion archive contains this article in editable \LaTeX\ and
PDF, the exact arithmetic replay, typed reference fixtures, integrity
checks, and documented clean-build and deterministic-pack commands.
It also includes the unchanged final Report 135 companion and its
Report 134 materials. The latter retain their own verifiers and seals;
the new replay checks their byte integrity before relying on them.
No third-party paper PDFs are redistributed.

The new finite checks cover the single-log coefficient identities and
squared-log coefficient products, triangular singular-vector matching,
the Gaussian coefficients through $\alpha_3$, the two exact kernel
pairings and $\kappa_4$, the finite $r_3$ shift weights, and the first
inverse-polynomial cancellations. Additional regularization checks use
explicit synthetic scalar sequences with exactly known model tails.
They check the implementation of subtraction and convergent weighted
sums. They are not numerical evaluations or error certificates for the
actual A217057 constants.

The verifier rejects undeclared or ill-typed fixture data and checks the
sealed source inventory. Normal and optimized Python runs are required
to produce the same result; the tests do not rely on removable
\texttt{assert} statements. Separate output locations and explicit
path guards protect the frozen package. Fresh clean builds and
repacking are checked for deterministic output. These finite
computations support the displayed algebra and reproducibility.
The uniform estimates, infinite-series convergence, and all-orders
claims are established by the proofs in this article, not by testing
finitely many coefficients.

\section{Further questions}
The all-orders statement, existence of the finite third constant,
convergent regularization algorithm, next exact logarithmic coefficient,
and arbitrary fixed-order inverse sandwich are proved here from the
exact inputs of Report 135. The following questions remain separate.
\begin{enumerate}
\item Can the regularized series be evaluated with explicit tail bounds
to certify useful decimal enclosures for $R,S,S_2,r_3$ and later
nonlogarithmic constants? Absolute convergence alone supplies no
numerical certificate at a stated truncation.
\item Can the local and off-chart estimates be made effective enough
to give practical constants and onset thresholds in the direct and
inverse remainders? A shrinking asymptotic error interval still does
not resolve exact integer rounding if it crosses an integer.
\item What is the growth of the coefficients as their order increases?
Convergence, Gevrey bounds, or summability of the infinite formal series
would require additional analysis. No exponentially small sector
expansion is established here.
\item Can the exact rational-diagonal description from Report 135 be
turned into a usable differential equation or recurrence that supports
certified large-index computation of the regularized constants?
\item Which other exact-pattern decompositions admit a comparable
weighted Banach expansion and a bilinear convolution reduction? The
coefficient theorem applies whenever its explicit hypotheses hold;
verifying those hypotheses for another pattern is a new problem.
\end{enumerate}

\begin{thebibliography}{9}
\bibitem{R134} Report 134, \emph{Permutations with exactly one increasing
subsequence of length four}, October 2, 2026. Unchanged \LaTeX, PDF,
exact-count code, and certificate included in the Report 135 companion.
\bibitem{R135} Report 135, \emph{Rational structure and logarithmic
asymptotics for A217057}, October 2, 2026. Final version with the
second correction and first logarithmic term; included unchanged with
its full reproducibility companion.
\bibitem{Isserlis} L. Isserlis, On a formula for the product-moment
coefficient of any order of a normal frequency distribution in any
number of variables, \emph{Biometrika} \textbf{12} (1918), 134--139.
\href{https://doi.org/10.1093/biomet/12.1-2.134}{doi:10.1093/biomet/12.1-2.134}.
\bibitem{DLMFharm} NIST Digital Library of Mathematical Functions,
\href{https://dlmf.nist.gov/5.4.E14}{equation 5.4.14}
(harmonic numbers and the digamma function) and
\href{https://dlmf.nist.gov/5.11.E2}{equation 5.11.2}
(fixed-order digamma asymptotics). Consulted October 2, 2026.
\bibitem{DLMFEM} NIST Digital Library of Mathematical Functions,
\href{https://dlmf.nist.gov/2.10.E1}{equation 2.10.1}
(Euler--Maclaurin with integral remainder) and
\href{https://dlmf.nist.gov/2.10.E8}{equation 2.10.8}
(harmonic asymptotics and the stated remainder bound).
Consulted October 2, 2026.
\end{thebibliography}
\end{document}
